% Submitted to The Annals of Statistics on ...
% Revised on ...
\documentclass[a4paper,12pt]{article}
%\documentclass[a4paper,titlepage]{article}

\usepackage{amsmath}
\usepackage{amssymb}
%\usepackage[dvips]{graphicx}
\usepackage{graphicx}
\usepackage[OT4]{fontenc}
%\usepackage[T1]{fontenc}
\usepackage[utf8]{inputenc}
\usepackage{tikz}
\usepackage{url}
\usepackage[numbers]{natbib}
\usepackage{bm}
\emergencystretch=20pt
\allowdisplaybreaks

\usepackage{setspace}
%\doublespacings

\newcommand{\abs}[1]{\left| #1 \right|}
\newcommand{\norm}[1]{\left|\mkern-1mu\left| #1 \right|\mkern-1mu\right|}
\newcommand{\aabs}[1]{\left\| #1 \right\|}
\newcommand{\okra}[1]{\left( #1 \right)}
\newcommand{\kwad}[1]{\left[ #1 \right]}
\newcommand{\klam}[1]{\left\{ #1 \right\}}
\newcommand{\ceil}[1]{\left\lceil #1 \right\rceil}
\newcommand{\floor}[1]{\left\lfloor #1 \right\rfloor}
\newcommand{\dzi}[1]{\left\langle #1 \right\rangle}
\newcommand{\ind}{\mathrel{\perp\mkern-10mu\perp}}
\DeclareMathOperator*{\argmin}{arg\, min}
\DeclareMathOperator*{\argmax}{arg\, max}
\DeclareMathOperator*{\hilberg}{hilb}
%\DeclareMathOperator{\atan}{atan}
\DeclareMathOperator{\essinf}{ess\, inf}
\DeclareMathOperator{\esssup}{ess\, sup}
\DeclareMathOperator{\card}{card}
\DeclareMathOperator{\supp}{supp}
\DeclareMathOperator{\cov}{Cov}
\DeclareMathOperator{\var}{Var}
\DeclareMathOperator{\mean}{\mathbf{E}}
\DeclareMathOperator{\median}{\mathbf{M}}
\newcommand{\peq}{\stackrel{+}{=}}
\newcommand{\pge}{\stackrel{+}{>}}
\newcommand{\ple}{\stackrel{+}{<}}
\newcommand{\boole}[1]{{\bf 1}{\klam{#1}}}

\newtheorem{conjecture}{Conjecture}
\newtheorem{definition}{Definition}
\newtheorem{example}{Example}
\newtheorem{corollary}{Corollary}
\newtheorem{proposition}{Proposition}
\newtheorem{theorem}{Theorem}
\newtheorem{lemma}{Lemma}
\newenvironment*{proof}{\begin{trivlist}\item[]
\noindent\textbf{Proof:}}{$\Box$\par\end{trivlist}}
\newenvironment*{proof*}[1]{\begin{trivlist}\item[]
\noindent\textbf{Proof of #1:}}{$\Box$\par\end{trivlist}}
%\renewcommand{\theenumi}{\roman{enumi}}
%\renewcommand{\labelenumi}{(\theenumi)}
%\renewcommand{\thesection}{\Roman{section}}
%\renewcommand{\thesubsection}{\Roman{section}-\Alph{subsection}}


\newcommand{\bibnamefont}[1]{\MakeUppercase{#1}}


\author{{\L}ukasz D\k{e}bowski\thanks{
    {\L}. D\k{e}bowski is with
    the Institute of Computer Science, Polish Academy of Sciences, 
    ul. Jana Kazimierza 5, 01-248 Warszawa, Poland 
    (e-mail: ldebowsk@ipipan.waw.pl).
  }
  % \thanks{Manuscript received ... ; revised ... .}
  %\thanks{}  
}


\title{Repetition and recurrence times in the case of \\ a stretched
  exponential growth} \date{}

\begin{document}

%\pagestyle{empty}   
\begin{titlepage}
%\titlepage
\maketitle

\begin{abstract}
  By an analogy to the duality between the recurrence time and the
  longest match length, we introduce a quantity dual to the maximal
  repetition length, which we call the repetition time.  Extending
  prior results, we sandwich the repetition time in terms of
  unconditional and conditional min-entropies. The condition for the
  upper bound resembles short memory in the sense developed in time
  series analysis. Our reasonings make a repeated use of dualities
  between so called times and so called counts that generalize the
  duality of the recurrence time and the longest match length.  We
  also discuss the analogy of these results with the
  Wyner-Ziv/Ornstein-Weiss theorem, which sandwiches the recurrence
  time in terms of Shannon entropies. We formulate the respective
  sandwich bounds in a way that applies to some cases of a stretched
  exponential growth. Such a growth was observed empirically for
  natural language. However, our results are too weak to apply to
  natural language directly because of its long-memory properties.
  % ca. 160 words
% \end{abstract}
%
% \begin{keywords}
  \\[1em]
  \textbf{Keywords:} stationary processes, recurrence time, maximal
  repetition, R\'enyi entropies, short memory, time-count duality
  % \end{keywords}
  \\[1em]
  \textbf{MSC2020:} 60G10, 94A17
\end{abstract}

% %MSC 2020: 
% % 94A29 - Source coding
% % 60G10 - Stationary processes
% % 94A17 - Measures of information, entropy


\end{titlepage}
% \pagestyle{plain}   

% \tableofcontents

\section{Introduction}
\label{secIntroduction}

Using appropriate concepts is often a gate to the simple theory. Some
instance of this rule will be exhibited in this note. The paper
concerns stochastic bounds for the maximal repetition length and its
newly introduced dual---called the repetition time.  Our development
will be analogous to the known duality between the recurrence time and
the longest match length \citep{WynerZiv89}.  Whereas the repetition
time is defined in this paper, there is a long history of
investigation of the recurrence time
\citep{Kac47,WynerZiv89,OrnsteinWeiss93,Kontoyiannis98,Ko12}, the
longest match length \citep{WynerZiv89,OrnsteinWeiss93,Szpankowski93,
  KontoyiannisOthers98,GaoKontoyiannisBienenstock08}, and the maximal
repetition length
\citep{DeLuca99,KolpakovKucherov99a,KolpakovKucherov99,CrochemoreIlie08,
  ErdosRenyi70,ArratiaWaterman89,Shields92b,Shields97,
  Szpankowski93,Szpankowski93a,KontoyiannisSuhov94,
  Debowski11b,Debowski17,Debowski18b}. These concepts appear in
various contexts in ergodic theory, information theory, theoretical
computer science, and quantitative linguistics, as we have indicated
by the references and we will detail further.  Establishing the
asymptotic behavior of these random variables for stationary ergodic
processes is of a practical importance since when we are confronted
with empirical data, such results can indicate what kind of a process
may generate these observations.

To be concrete, let $(X_i)_{i\in\mathbb{Z}}$ be an arbitrary
stochastic process over a countable alphabet. The well-researched
recurrence time $R^{(1)}_k$ is the first position in the future on
which a copy of string $X_1^k$ reappears,
\begin{align}
  R^{(1)}_k:=\inf\klam{\infty}\cup\klam{i\ge 1: X_{i+1}^{i+k}=X_1^k}.
\end{align}
By contrast, we define here the repetition time $R^{(2)}_k$ as the
first position in the future on which a copy of any previously seen
string $X_{j+1}^{j+k}$ occurs,
\begin{align}
  R^{(2)}_k:=\inf\klam{\infty}\cup\klam{i\ge 1:
  X_{i+1}^{i+k}=X_{j+1}^{j+k}\text{ for some } 0\le j<i}.
\end{align}
If $(X_i)_{i\in\mathbb{Z}}$ is a process over a $D$-ary alphabet then
$R^{(2)}_k\le D^k$. In any case, $R^{(2)}_k\le R^{(1)}_k$. Precisely,
we have
\begin{align}
  \label{Minimum}
  R^{(2)}_k=\inf_{i\ge 0}\okra{i+R^{(1)}_k\circ T^i}
\end{align}
for the shift automorphism $X_i\circ T=X_{i+1}$, since observing a
repetition is equivalent to observing a recurrence of some particular
string.

In parallel, we can discuss variables that are dual to random
variables $R^{(\gamma)}_k$. The longest match length $L^{(1)}_n$ is
\begin{align}
  L^{(1)}_n:=\max\klam{0}\cup\klam{k\ge 1: X_{i+1}^{i+k}=X_1^k
  \text{ for some } 0< i\le n-k},
\end{align}
whereas the maximal repetition length $L^{(2)}_n$ is
\begin{align}
  L^{(2)}_n:=\max\klam{0}\cup\klam{k\ge 1: X_{i+1}^{i+k}=X_{j+1}^{j+k}
  \text{ for some } 0\le j< i\le n-k}.
\end{align}
Both lengths $L^{(\gamma)}_n$ are functions of prefix $X_1^n$ and
$n\ge L^{(2)}_n\ge L^{(1)}_n$. For $\gamma\in\klam{1,2}$, we have the
dualities
\begin{align}
  \label{DualityOne}
  R^{(\gamma)}_k&=\inf\klam{\infty}\cup\klam{n\ge 1: L^{(\gamma)}_{n+k}\ge k},
  \\
  \label{DualityTwo}
  L^{(\gamma)}_n&=\max\klam{0}\cup\klam{k\ge 1: R^{(\gamma)}_k\le n-k}.
\end{align}
Hence $R^{(\gamma)}_k>n\iff L^{(\gamma)}_{n+k}<k$
\citep{WynerZiv89}. Thus the sequence of the dual random variables
carries the same information as the sequence of the primal ones. Both
$R^{(\gamma)}_k$ and $L^{(\gamma)}_n$ are non-decreasing functions of
$k$ and $n$, respectively.

These dualities are an instance of a more general duality that we
propose to call the time-count duality.
\begin{theorem}[time-count duality, cf.\ \citep{WynerZiv89}]
  \label{theoTimeCount}
  % Denote $\mathbb{N}_+:=\mathbb{N}\cup\klam{0,\infty}$
  Suppose that function
  $F:\mathbb{N}\rightarrow\mathbb{N}\cup\klam{0}$ (called the count)
  is non-decreasing. Define the dual function
  $C:\mathbb{N}\rightarrow\mathbb{N}\cup\klam{\infty}$ (called
  the time) as
  \begin{align}
    \label{Time}
    C(m)&:=\inf\klam{\infty}\cup\klam{n\ge 1: F(n)\ge m}.
  \end{align}
  We claim that:
  \begin{enumerate}
  \item $C(m)>n\iff F(n)<m$;    
  \item $F(n)=\max\klam{0}\cup\klam{m\ge 1: C(m)\le n}$;    
  \item Function $C$ is non-decreasing.
  \end{enumerate}
\end{theorem}
\emph{Remark:} Since time $C$ is non-decreasing, we may apply
construction (\ref{Time}) to $C$ rather than $F$. But then we obtain a
shifted function
\begin{align}
  F'(n)
  &:=\inf\klam{\infty}\cup\klam{m\ge 1: C(m)\ge n}
    \nonumber\\
  &=
  \inf\klam{\infty}\cup\klam{m\ge 1: C(m)>n-1}
    \nonumber\\
  &=
  \inf\klam{\infty}\cup\klam{m\ge 1: F(n-1)<m}
    \nonumber\\
  &=
    \inf\klam{\infty}\cup\klam{m\ge 1: F(n-1)\le m-1}
    \nonumber\\
  &=F(n-1)+1.
\end{align}
\begin{proof}
  Since $F$ is non-decreasing, condition $F(n)<m$ implies $C(m)>n$.
  On the other hand, by definition, if $F(n)\ge m$ then $C(m)\le n$.
  Thus we have equivalence $C(m)>n\iff F(n)<m$.  In view of this
  equivalence, we may write
  \begin{align}
    F(n)&=\max\klam{m\ge 0: F(n)\ge m}
    \nonumber\\
        &=\max\klam{0}\cup\klam{m\ge 1: C(m)\le n}
    .
  \end{align}
  Suppose that $m_1\le m_2$. We have
  \begin{align}
    \klam{n\ge 1: F(n)\ge m_1}\supset \klam{n\ge 1: F(n)\ge
    m_2}.
  \end{align}
  Since $F$ is non-decreasing, both sets are intervals. Thus
  $C(m_1)\le C(m_2)$.
\end{proof}

Easily, the time-count duality carries over to some extreme values.
\begin{theorem}
  \label{theoExtreme}
  Suppose that functions $F^B$ and $C^B$ are dual for all
  $B\in\beta$. Then functions $F:=\sup_{B\in\beta} F^B$ and
  $C:=\inf_{B\in\beta} C^B$ are also dual.
\end{theorem}
\begin{proof}
  The claim follows by equivalence
  \begin{align}
    C(m)>n
    &\iff\forall_{B\in\beta}\,C^B(m)>n
      \nonumber\\
    &\iff
      \forall_{B\in\beta}\,F^B(n)<m
      \iff F(n)<m.
  \end{align}
  Hence
  \begin{align}
    C(m)&=\inf\klam{\infty}\cup\klam{n\ge 1: C(m)\le n}
          \nonumber\\
        &=\inf\klam{\infty}\cup\klam{n\ge 1: F(n)\ge m}.
  \end{align}
\end{proof}
For a discrete stochastic process $(X_i)_{i\in\mathbb{Z}}$ and a
certain fixed value $x$, some basic time-count duality arises for
counts $F_n(x):=\sum_{j=1}^n \boole{X_j=x}$ and times
$C_m(x):=\inf\klam{\infty}\cup\klam{n\ge 1: F_n(x)\ge m}$.  We will
make use of such dualities in this paper and we will link them with
the recurrence time via Theorem \ref{theoExtreme}.

Speaking of the recurrence and repetition time again, let us comment
on the received theoretical results in the stochastic setting. We
begin with the recurrence time and the longest match length.  Let
$\log x$ denote the natural logarithm of $x$.  Consider a discrete
random variable $X$. The Shannon entropy is
$H_1(X):=\mean\kwad{-\log P(X)}$. Let $(X_i)_{i\in\mathbb{Z}}$ be a
stationary ergodic process over a finite alphabet. The celebrated
Wyner-Ziv/Ornstein-Weiss (WZ/OW) theorem
\citep{WynerZiv89,OrnsteinWeiss93} asserts
\begin{align}
  \label{WZOW}
  \lim_{k\to\infty} \frac{\log n}{L^{(1)}_n}
  =\lim_{k\to\infty} \frac{1}{k}\log R^{(1)}_k
  =h_1:=\lim_{k\to\infty} \frac{H_1(X_1^k)}{k}
  \text{ a.s.}
\end{align}
Thus the rates of the recurrence time and the longest match length are
essentially dictated by the rate of the Shannon entropy.  The Shannon
entropy rate $h_1$ exists for any stationary process.

Let us proceed to the maximal repetition length and the repetition
time.  For memoryless sources, the maximal repetition length
$L^{(2)}_n$ grows at a similar rate as the longest match length
$L^{(1)}_n$ \citep{ErdosRenyi70,ArratiaWaterman89}. Probably for this
reason, the problem that $L^{(2)}_n$ can grow much faster than
$L^{(1)}_n$ for other stationary processes was first recognized as
late as by \citet{Shields92b,Shields97}. Define the collision entropy
as $H_2(X):=-\log\mean P(X)$, where $H_2(X)\le H_1(X)$ by the Jensen
inequality.  A technically complicated result by
\citet{Szpankowski93,Szpankowski93a}, involving suffix trees, can be
rewritten as
\begin{align}
  \label{Szpankowski}
  \lim_{k\to\infty} \frac{\log n}{L^{(2)}_n}
  =\lim_{k\to\infty} \frac{1}{k}\log R^{(2)}_k
  =h_2:=\lim_{k\to\infty} \frac{H_2(X_1^k)}{k}
  \text{ a.s.}
\end{align}
provided we have a uniform bound on the pointwise mutual information
\begin{align}
  \label{Kfinite}
  \sup_{n,m\ge 1}\esssup
     \frac{P(X_1^{n+m})}{P(X_1^n)P(X_{n+1}^{n+m})}
     < \infty.
\end{align}
Quantity $h_2$ is called the collision entropy rate. Whereas limit
$h_1$ exists for any stationary process, limit $h_2$ exists if
(\ref{Kfinite}) holds additionally. The collision entropy rate can be
substantially smaller than the Shannon entropy rate.

Condition (\ref{Kfinite}) is satisfied by finite-state processes
\citep[Theorem 11.3]{Debowski21} but it does not cover all important
cases---such as stationary sources that do not satisfy the finite
energy condition \citep{Shields92b,Shields97}, do not meet the
Doeblin condition \citep[Theorem 11.7]{Debowski21}, or are not diluted
by random noise \citep{Shields97}, \citep[Theorem 11.6]{Debowski21}.
We became interested in characterizing the maximal repetition length
because of our long-standing interests in statistical language
modeling \citep{Debowski21}.  In fact, for natural language,
experiments suggest that the Shannon entropy rate is positive,
$h_1>0$, \citep{Shannon51,CoverKing78,BehrOthers03} but the mutual
information probably grows like a power-law, called Hilberg's law
\citep{Hilberg90, TakahiraOthers16, KaplanOther20, Debowski21}, which
contradicts condition (\ref{Kfinite}). Moreover, as discovered
empirically in \citep{Debowski12b,Debowski15f} for novels in English,
German, and French, we have a power-law logarithmic growth of the
maximal repetition length
\begin{align}
  L^{(2)}_n\propto \log^\alpha n,\quad \alpha\approx 3.
\end{align}
This is equivalent to the stretched exponential growth of the
repetition time
\begin{align}
  \log R^{(2)}_k\propto k^\beta,\quad \beta\approx 1/3.
\end{align}

The goal of this paper is to analyze this case of a stretched
exponential growth, strengthening and beautifying our earlier
theoretical accounts and cumbersome statements
\citep{Debowski15f,Debowski18b,Debowski21}. We will do it by using the
minimal times $R^{(\gamma)}_k$ rather than the maximal lengths
$L^{(\gamma)}_n$. We will have three main results of a varying degree
of originality.

A few explanations regarding our framework are in place.  First, the
Hilberg exponent of a sequence $(a_k)_{k\in\mathbb{N}}$ is defined as
\begin{align}
  \label{Exponent}
  \hilberg_{k\rightarrow\infty} a_k
  :=\limsup_{k\rightarrow\infty} \max\klam{0,\frac{\log
  a_k}{\log k}}
  .
\end{align}
It measures the rate of asymptotic power-law growth,
$\hilberg_{k\rightarrow\infty} k^\beta=\max\klam{0,\beta}$. Some
properties of the Hilberg exponent are summarized in \citep[Section
8.1]{Debowski21}. We write succinctly
\begin{align}
  a_k\lesssim b_k\iff
  \hilberg_{k\rightarrow\infty} a_k\le
  \hilberg_{k\rightarrow\infty} b_k,
  \\
  a_k\sim b_k\iff
  \hilberg_{k\rightarrow\infty} a_k=
  \hilberg_{k\rightarrow\infty} b_k.
\end{align}
The Hilberg exponent is useful for characterizing the power-law growth
of mutual information, i.e., Hilberg's law \citep[Chapter
8]{Debowski21}. Hence comes its name.

Second, for a parameter $\gamma\in(0,1)\cup(1,\infty)$, the
R\'enyi-Arimoto conditional entropy of a discrete random variable $X$
given a variable $Y$ is defined as
\begin{align}
  \label{RenyiArimoto}
  H_\gamma(X|Y):=
  -\frac{\gamma}{\gamma-1}\log
  \mean\kwad{\sum_x p(X=x|Y)^{\gamma}}^{1/\gamma}
  ,
\end{align}
whereas for $\gamma\in\klam{0,1,\infty}$, we define
\begin{align}
  H_\gamma(X|Y):=\lim_{\delta\rightarrow\gamma} H_\delta(X|Y).
\end{align}
We also put $H_\gamma(X):=H_\gamma(X|1)$. The special cases are
called:
\begin{itemize}
\item the Hartley entropy
$H_0(X|Y)= \log \esssup\card\klam{x: P(X=x|Y)>0}$,
\item  the Shannon
entropy $H_1(X|Y)=\mean\kwad{-\log P(X|Y)}$, 
\item the collision entropy
$H_2(X|Y)=-2\log \mean\kwad{\sum_{x\in\mathbb{X}} P(X=x|Y)^2}^{1/2}$,
\item the min-entropy $H_\infty(X|Y)=-\log \mean\max_{x} P(X=x|Y)$.
\end{itemize}
We have $H_\gamma(X)\ge H_\delta(X)$ for $\gamma<\delta$, whereas
$H_\gamma(X)\le \frac{\gamma}{\gamma-1}H_\infty(X)$ for $\gamma>1$.
Thus all R\'enyi-Arimoto entropies with $\gamma>1$ are of the same
magnitude.  Moreover, we have a generalized chain rule
\begin{align}
  \label{Chain}
  H_\gamma(X|Y,Z)\le H_\gamma(X|Y)\le H_\gamma(U,X|Y)\le
  H_0(U|Y)+H_\gamma(X|Y,U).
\end{align}
A survey of the R\'enyi-Arimoto entropies can be found in
\citep{FehrBerens14} and also in \citep[Section 9.1]{Debowski21}.

Let us come back to the recurrence and repetition times.  In the body
of the paper, we will prove the following Theorems
\ref{theoMainOne}--\ref{theoMainThree}, which bound the quantities of
interest comprehensively but a noticeable novelty can be attributed
only to Theorem \ref{theoMainTwo}.
\begin{theorem}
  \label{theoMainOne}
  Consider a stationary ergodic process $(X_i)_{i\in\mathbb{Z}}$ over
  a $D$-ary alphabet. Define $N_k:=\max\klam{N: N\log N\le D^k}$. We
  have a sandwich bound
  \begin{align}
    % h_1 k\sim
    \log P(X_1^k|X_{k+1}^{N_k})\lesssim
    \log R^{(1)}_k\lesssim
    -\log P(X_1^k)\lesssim
    H_1(X_1^k)
    \text{ a.s.},
  \end{align}
  where the right-most inequality $\lesssim$ becomes the equivalence
  $\sim$ for
  \begin{align}
    \limsup_{k\to\infty} \frac{\sqrt{\var\log P(X_1^k)}}{H_1(X_1^k)}<1.
  \end{align}
\end{theorem}
\emph{Remark:} Theorem \ref{theoMainOne} applies the same sandwich
bound by \citet{Kontoyiannis98} as the simple proof of the WZ/OW
theorem (\ref{WZOW}) presented in \citep[Theorem 9.12]{Debowski21}.
We also recall that $h_1=\mean\kwad{-\log P(X_1|X_2^\infty)}$ for a
stationary process. Consequently if $h_1=0$ then
$P(X_1^k|X_{k+1}^\infty)=1$ almost surely. Quantity
$\var\log P(X_1^k)$ is called varentropy \citep{KontoyiannisVerdu13}.

For the next theorem, we will apply some concepts of short memory that
resemble the notion generally adopted in time series analysis
\cite{Beran94}, namely, summability of autocorrelations.  For a
stationary process $(X_i)_{i\in\mathbb{Z}}$, we define random counts
\begin{align}
  F_n(x_1^k)&:=\sum_{i=0}^{n-1}\boole{X_{i+1}^{i+k}=x_1^k}.
\end{align}
We consider the following notions.
\begin{definition}[short memory]
  \mbox{}
  \begin{enumerate}
  \item Process $(X_i)_{i\in\mathbb{Z}}$ is called short
    memory (I) if $\log \gamma_k\sim 0$ for coeffcients
    \begin{align}
      \label{Gamma}
      \gamma_k&:=\sup_{n\in\mathbb{N}}\max_{x_1^k}
                \frac{\var F_n(x_1^k)}{\mean F_n(x_1^k)}.
    \end{align}
  \item Process $(X_i)_{i\in\mathbb{Z}}$ is called short memory (II)
    if $\log \delta_k\sim 0$ for coeffcients
    \begin{align}
    \label{Delta}
      \delta_k:=
      \sup_{n\in\mathbb{N}}\frac{\mean\max_{x_1^k}
      F_n(x_1^k)}{\max_{x_1^k}\mean F_n(x_1^k)}.
    \end{align}
  \end{enumerate}
\end{definition}
\emph{Remark:} Short memory (I) is stronger than summability of
autocorrelations.  Namely, we require that the sums of correlations
for variables $\boole{X_{i+1}^{i+k}=x_1^k}$ are finite for each block
length $k$ and they do not grow too fast with respect to $k$.  For
example, for an IID process $(X_i)_{i\in\mathbb{Z}}$, we have
$\gamma_k=O(k)$. As for the condition of short memory (II), if counts
$F_n(x_1^k)$ were IID for a fixed $n$ and $x_1^k$ varying over $D$-ary
strings of length $k$ then it would be reasonable to expect
$\delta_k=O(k)$ by the extreme value theory \citep{HaanFereira06}.  In
view of the asymptotic equipartition, i.e., the
Shannon-McMillan-Breiman (SMB) theorem \citep{Shannon48, Breiman57,
  Chung61, Barron85, AlgoetCover88}, we suppose that for many
stationary processes, counts $F_n(x_1^k)$ are close to IID for large
$n$ so $\log \delta_k\sim 0$ seems a reasonable condition.

Consequently, we may demonstrate these facts.
\begin{theorem}
  \label{theoMainTwo}
  \mbox{}
  \begin{enumerate}
  \item Consider a stationary short memory (I) process
    $(X_i)_{i\in\mathbb{Z}}$ over a finite alphabet. We have an upper
    bound
  \begin{align}
    \label{MainTwoOne}
    \log R^{(2)}_k\lesssim
    H_\infty(X_1^k)
    \text{ a.s.}
  \end{align}
  \item Consider a stationary short memory (II) process
    $(X_i)_{i\in\mathbb{Z}}$ over a finite alphabet. We have a lower
    bound
  \begin{align}
    \label{MainTwoTwo}
    H_\infty(X_1^k)\lesssim \log R^{(2)}_k
    \text{ a.s.}
  \end{align}
  \end{enumerate}
\end{theorem}
\emph{Remark:} In \citep{Debowski18b,Debowski21}, a trivial bound in
terms of Shannon entropy $H_1(X_1^k)$ stemming from inequality
$L^{(2)}_n\ge L^{(1)}_n$ was discussed instead of the upper bound
(\ref{MainTwoOne}). In the previous preprint version of this article,
we claimed erroneously that $\log R^{(2)}_k$ can be generally upper
bounded by $H_\infty(X_1^k)$. The phenomenon of burstiness, described
in Section \ref{secRepetition} likely excludes such a possibility.
Bounds (\ref{MainTwoOne})--(\ref{MainTwoTwo}) pertain to a class of
short memory (I)--(II) processes, which is somewhat comparable to
processes satisfying condition (\ref{Kfinite}), examined in
\citep{Szpankowski93,Szpankowski93a}. However, the conditions of short
memory (I)--(II) seem easier to verify empirically than condition
(\ref{Kfinite}). Of course, under short memory (I) and (II) combined,
we have $\log R^{(2)}_k\sim H_\infty(X_1^k)$ almost surely.

Unfortunately, although the class of short memory (I)--(II) processes
may be larger than one defined by condition (\ref{Kfinite}), it is
still too restrictive for our purposes due to Taylor's law. Taylor's
law is an empirical power law for scaling of variance in natural
language
\citep{GerlachAltmann14,KobayashiTanakaIshii2018,TanakaIshii21} that
implies $\gamma_k=\infty$ in particular. This hinders the application
of bound (\ref{MainTwoOne}) to natural language. Thus we are not fully
satisfied with Theorem \ref{theoMainTwo}.

In passing, we notice also a general lower bound for all stationary
ergodic processes. This theorem somewhat strengthens \citep[Theorem
8]{Debowski18b}.
\begin{theorem}
  \label{theoMainThree} 
  Consider a stationary process $(X_i)_{i\in\mathbb{Z}}$ over a
  countable alphabet. Define the context length
    \begin{align}
      \label{ContextLength}
      I_k&:=\min\klam{i:\log i\ge H_\infty(X_1^k|X_{k+1}^{k+i})}.
    \end{align}
    We have a lower bound
    \begin{align}
      \label{MainThree}
      H_\infty(X_1^k|X_{k+1}^{k+I_k})\lesssim
      \log I_k\lesssim
      \log R^{(2)}_k
      \text{ a.s.}
      ,
    \end{align}
    where the left-most inequality $\lesssim$ becomes the equivalence
    $\sim$ for a finite alphabet.  
\end{theorem}
\emph{Remark:} As for the lower bound (\ref{MainThree}), we may
substitute
\begin{align*}
  I_k\to\min\klam{i:\frac{\gamma}{\gamma-1}\log i\ge
  H_\gamma(X_1^k|X_{k+1}^{k+i})}
\end{align*}
for any $\gamma>1$.  Context length $I_k$ can be much shorter than the
pessimistic context length $D^k$ for a $D$-ary alphabet that was
proved sufficient in \citep[Theorem 8]{Debowski18b}.
% Moreover, bounding $R^{(2)}_k$ rather than
% $L^{(2)}_n$ leads to more elegant statements than produced in works
% \citep{Debowski18b,Debowski21}.

This paper intentionally covers theorems rather than examples.
Construction of instances of stationary ergodic processes that satisfy
proportionality $L^{(2)}_n\propto \log^\alpha n$ with exponent
$\alpha>1$ and the Shannon entropy rate $h_1>0$, stated as an open
problem in book \citep{Debowski21}, is relegated to another
article. The essential idea, inspired by Theorem \ref{theoMainTwo}, is
not so difficult---the source should emit sufficiently long words in a
memoryless fashion. However, because of other involved linguistic
concepts, such as Hilberg's law \citep{Hilberg90, TakahiraOthers16,
  KaplanOther20, Debowski21}, and unfinished calculations, we prefer
to publish this detailed result elsewhere.

The organization of this note is as follows. In Section
\ref{secRecurrence}, we recall the known sandwich bound for the
recurrence time. In Section \ref{secRepetition}, we prove a new
sandwich bound for the repetition time under short memory
(I)--(II). In Section \ref{secYetAnother}, we discuss a lower bound for
the repetition time under plain stationarity.  Concluding, in Section
\ref{secFinal}, we specialize these results to the Hilberg exponents.

% -- How much can different entropies differ?, Szpankowski's results
% on $L^{(\gamma)}_n$ for $\gamma\ge 3$ suggest that entropy rates
% $h_\gamma$ are equal.


\section{Sandwiching recurrence times}
\label{secRecurrence}

There is a celebrated lemma by \citeauthor{Kac47}. Namely, the
conditional expectation of the recurrence time is the inverse
probability of the initial string.
\begin{lemma}[\citep{Kac47}]
  \label{theoKac}
  For a stationary ergodic process $(X_i)_{i\in\mathbb{Z}}$ over a
  countable alphabet,
  \begin{align}
    \mean\okra{R^{(1)}_k\middle|X_1^k}
    =
    \frac{1}{P(X_1^k)}.
  \end{align}
\end{lemma}
There is also a converse result by \citet{Kontoyiannis98}, which was
strengthened by \citet{Debowski18b} from conditioning on
$X_{k+1}^\infty$ to conditioning on $X_{k+1}^{D^k}$. Here we shorten
the conditioning to $X_{k+1}^{N}$, where $N<D^k$.
\begin{lemma}[cf.\ \citep{Kontoyiannis98}, \mbox{\citep[Lemma 2]{Debowski18b}}]
  \label{theoKontoyiannis}
  For a process $(X_i)_{i\in\mathbb{Z}}$ over a $D$-ary alphabet and
  an integer constant $N$ where $1\le N\le D^k$,
  \begin{align}
    \mean\okra{\frac{1}{\min\klam{R^{(1)}_k,N}
    P(X_1^k|X_{k+1}^N)}\middle|X_{k+1}^N}
    \le
    1+\log N+\frac{D^k-N}{N}.
  \end{align}
\end{lemma}
\begin{proof}
  Define the waiting time for string $x_1^k$ as
  \begin{align}
   R(x_1^k):=\inf\klam{\infty}\cup\klam{i\ge 1: X_{i+1}^{i+k}=x_1^k}. 
  \end{align}
  We have
  \begin{align}
    \mean\okra{\frac{1}{\min\klam{R^{(1)}_k,N}
    P(X_1^k|X_{k+1}^N)}\middle|X_{k+1}^N}
    =
    \sum_{x_1^k}\frac{1}{\min\klam{R(x_1^k),N}}.
  \end{align}
  We observe that for each $i\ge 1$ there is at most one string
  $x_1^k$ such that $R(x_1^k)=i$. Hence, we have a uniform bound
\begin{align}
  \sum_{x_1^k}\frac{1}{\min\klam{R(x_1^k),N}}
  &\le
  \sum_{i=1}^{D^k}\frac{1}{\min\klam{i,N}}
  \le
    1+\int_{1}^{N}\frac{1}{u}du+\int_{N}^{D^k}\frac{1}{N}du
    \nonumber\\
  &\le
  1+\log N+\frac{D^k-N}{N}.
\end{align}
\end{proof}

The following sandwich bound essentially appears in the simple proof
of the WZ/OW theorem in \citep[Theorem 9.12]{Debowski21}, which was
originally proved using other techniques, cf.\
\citep{WynerZiv89,OrnsteinWeiss93}. We recall our simpler derivation
for the analogy with Section \ref{secRepetition}. Applying Lemma
\ref{theoKontoyiannis}, we slightly shorten the length of conditioning
compared to the results of \citep{Debowski18b}.
\begin{proposition}
  \label{theoSandwichRecurrence}
  Consider a stationary ergodic process $(X_i)_{i\in\mathbb{Z}}$ over
  a $D$-ary alphabet. Define $N_k:=\max\klam{N: N\log N\le D^k}$.
  Assume that $\sum_{k=1}^\infty \rho_k<\infty$ for some $\rho_k>0$.
  Almost surely, for sufficiently large $k$ we have
\begin{align}
  \label{SandwichRecurrence}
  -\log P(X_1^k|X_{k+1}^{N_k})
  +\log \rho_k-\log k
  <
  \log R^{(1)}_k
  <
  -\log P(X_1^k)
  -\log \rho_k
  .
\end{align}
\end{proposition}
\begin{proof}
  As \citet{Kontoyiannis98} has shown, by applying the Markov
  inequality to Lemmas \ref{theoKac} and \ref{theoKontoyiannis}, we
  obtain two complementary bounds
\begin{align}
  \label{ProbRecurrenceI}
  P\okra{R^{(1)}_k\ge \frac{C}{P(X_1^k)}}
  &\le C^{-1}
  ,
  \\
  \label{ProbRecurrenceII}
  P\okra{R^{(1)}_k\le \frac{C}{P(X_1^k|X_{k+1}^{N_k})}}
  &\le C\okra{1+\log N_k+\frac{D^k-N_k}{N_k}}
  .
\end{align}
Setting $C=\rho_k^{-1}$ in (\ref{ProbRecurrenceI}) and $C=k^{-1}\rho_k$ in
(\ref{ProbRecurrenceII}) we obtain (\ref{SandwichRecurrence}) by the
Borel-Cantelli lemma.
\end{proof}

The WZ/OW theorem (\ref{WZOW}) follows by Proposition
\ref{theoSandwichRecurrence} combined with the
Shannon-McMillan-Breiman (SMB) theorem \citep{Shannon48, Breiman57,
  Chung61, Barron85, AlgoetCover88}, which asserts
\begin{align}
  \label{SMB}
  \lim_{k\to\infty}
  \frac{1}{k}\okra{-\log P(X_1^k)}=
  \lim_{k\to\infty}
  \frac{1}{k}\okra{-\log P(X_1^k|X_{k+1}^\infty)}=
  h_1 \text{ a.s.}
\end{align}

\section{Sandwiching repetition times}
\label{secRepetition}

There is a result by \citeauthor{ChenMoy59} that generalizes the Kac
lemma to successive recurrence times.  Let $T$ be a stationary
automorphism of the probability space.  Consider an event $B$ such
that $P(B)>0$. Define the random set
\begin{align}
  \mathcal{C}^B(\omega):=\klam{j\in\mathbb{Z}: T^{-j}\omega\in B}
\end{align}
If $P(B)>0$ then set $\mathcal{C}^B$ contains infinitely many negative
and positive integers almost surely, by the Poincar\'e recurrence
theorem. Then we can arrange the elements of set $\mathcal{C}^B$ into
the unique sorted sequence $(C^B_r)_{r\in\mathbb{Z}}$ defined by
conditions: $\bigcup_{r\in\mathbb{Z}}\klam{C^B_r}=\mathcal{C}^B$,
$C^B_r>C^B_{r-1}$ for $r\in\mathbb{Z}$, $C^B_1\ge 1$, and
$C^B_0\le 0$.  The successive recurrence times are defined as
\begin{align}
  \label{Successive}
  W^B_r:=C^B_r-C^B_{r-1}.
\end{align}
Observe that $(C^B_0=0)=B$. The result by Chen Moy is as follows.
\begin{lemma}[\citep{ChenMoy59}]
  \label{theoChenMoy}
  Process $(W^B_r)_{r\in\mathbb{N}}$ is stationary with respect to
  conditional measure $P(\cdot|B)$. Moreover, if automorphism $T$
  is ergodic then for $r\in\mathbb{N}$, we have
  \begin{align}
    \mean\okra{W^B_r\middle|B}
    =
    \frac{1}{P(B)}.
  \end{align}
\end{lemma}

Due to the phenomenon of burstiness, this result cannot be
strengthened to the unconditional expectation.
\begin{example}[periodic process with burstiness]
  Consider a periodic process $(X_i)_{i\in\mathbb{Z}}$ in which the
  pattern of $n_1$ ones followed by $n_0$ zeros repeats
  cyclically. Let $B=(X_0=1)$. Then $P(C^B_0=0)=P(B)=n_1/(n_0+n_1)$
  and $P(C^B_0=-i)=1/(n_0+n_1)$ for $i=1,...,n_0$. For $C_0<0$, we
  have $W^B_{kn_1+l}=1+n_0\boole{l=1}$ for
  $k\in\mathbb{N}\cup\klam{0}$ and $l=1,...,n_1$. Hence
\begin{align}
  \mean W^B_{kn_1+l}&=1+\frac{n_0+n_0^2\boole{l=1}}{n_0+n_1}.
\end{align}
For $n_1=1$, we have $\mean W^B_r=1+n_0=1/P(B)$. However, process
$(W^B_r)_{r\in\mathbb{N}}$ is non-stationary for other values of
$n_1$, which correspond to a burst of $1$'s followed a single lull
that is arbitrarily longer than $1/P(X_0=1)$.
\end{example}
We may expect that the phenomenon of burstiness is particularly acute
in statistical models of natural language
\citep{Katz96,AltmannPierrehumbertMotter09,Debowski15f}.

Let us try to work around this problem.  In the following, we denote
the number of visits to $B$ by time $n\in\mathbb{N}$ as
\begin{align}
  F^B_n(\omega):=\sum_{j=1}^n \boole{T^{-j}\omega\in B}.
\end{align}
We will call variables $F^B_n$ counts, whereas variables $C^B_m$ will
be called times since for $m,n\in\mathbb{N}$ we have the time-count
duality between times $C^B_m$ and counts $F^B_n$. Namely,
\begin{align}
  C^B_m&=\inf\klam{\infty}\cup\klam{n\ge 1: F^B_n\ge m},
  \\
  F^B_n&=\max\klam{0}\cup\klam{m\ge 1: C^B_m\le n}.
\end{align}
Thus we have equivalence $C^B_m>n\iff F^B_n<m$ by Theorem
\ref{theoTimeCount}.  What is interesting, although they carry the
same information, counts $F^B_n$ are much easier to analyze than times
$C^B_m$.  We have expectation $\mean F^B_n=nP(B)$ by stationarity,
whereas the Cauchy-Schwarz inequality yields further a bound for
variance
\begin{align}
 \var F^B_n\le n^2P(B)(1-P(B)).
\end{align}
Under short memory in the sense adopted in time series analysis
\citep{Beran94}, we have a stronger bound, namely,
$\var F^B_n\le\gamma nP(B)$ for a $\gamma<\infty$.

As a result, despite the phenomenon of burstiness, we can supply a
simple lower bound for the expectation of times $C^B_m$---by their
duality to counts $F^B_n$.
\begin{lemma}
  \label{theoUnconditional}
  If automorphism $T$ is stationary then for $m\in\mathbb{N}$, we have
  \begin{align}
    \mean C^B_m
    \ge
    \frac{m}{2P(B)}.
  \end{align}
\end{lemma}
\begin{proof}
  In view of equivalence $C^B_m>n\iff F^B_n<m$, we have
  $P(C^B_m>n)=P(F^B_n<m)$.  Hence, we may bound the expectation of
  $C^B_m$, using the Markov inequality as follows
  \begin{align}
    \mean C^B_m
    &=
    \int_0^\infty P(C^B_m>\floor{n}) dn
    =
    \int_0^\infty P(F^B_{\floor{n}}<m) dn
      \nonumber\\
    &=
      \int_0^\infty \kwad{1-P(F^B_{\floor{n}}\ge m)} dn
      \ge
      \int_0^\infty \kwad{1-\frac{\mean F^B_n}{m}}_+ dn
      \nonumber\\
    &\ge
      \int_0^\infty \kwad{1-\frac{nP(B)}{m}}_+ dn
      =
      \int_0^{m/P(B)} \kwad{1-\frac{nP(B)}{m}}_+ dn
      \nonumber\\
    &=
      \frac{m}{2P(B)}.
  \end{align}
\end{proof}

Assuming short memory in the sense of time series analysis
\citep{Beran94}, we may also prove an upper bound for the expectation
of variable $\min_{B\in\beta} C^B_m$, where $\beta$ is a finite
partition of the event space. Since variable $\min_{B\in\beta} C^B_m$
is upper bounded almost surely in this case then, in the following
proof, we can circumvent a problem with insummability of the harmonic
series and obtain a bound that is not void.
\begin{lemma}
  \label{theoShortMemory}
  Assume a stationary automorphism $T$ and a finite partition $\beta$.
  Suppose that
  \begin{align}
    \gamma:=
    \sup_{n\in\mathbb{N}}\max_{B\in\beta}\frac{\var F^B_n}{\mean F^B_n}<\infty.
  \end{align}
  Then for $m\in\mathbb{N}$, we have
  \begin{align}
    \mean \min_{B\in\beta} C^B_m\le
    \frac{1}{\max_{B\in\beta} P(B)}\okra{
    2m+\gamma\log\frac{(m-1)\card\beta+\gamma}{\gamma}
    }+1
    .
  \end{align}
\end{lemma}
\begin{proof}
  The Paley-Zygmund inequality for a random variable $Y\ge 0$ states
  that
  \begin{align}
    P(Y\le y)\le\frac{\var Y}{\var Y+(\mean Y-y)_+^2}.
  \end{align}
  Let $N\in\mathbb{N}$ be a constant. Then the duality yields
  \begin{align}
    \mean \min\klam{C^B_m,N}
    &=
    \int_0^{N} P(C^B_m>\floor{n}) dn
    =
      \int_0^{N} P(F^B_{\floor{n}}<m) dn
      \nonumber\\
    &\le \int_0^{N} \frac{\var F^B_{\floor{n}}
      dn}{\var F^B_{\floor{n}}+(\mean F^B_{\floor{n}}-m)_+^2}
      \nonumber\\
    & \le \int_0^{N} \frac{\gamma \floor{n}P(B)
      dn}{\gamma \floor{n}P(B)+(\floor{n}P(B)-m)_+^2}
      \nonumber\\
    &\le \int_0^{N}
      \frac{\gamma dn}{((n-1)P(B)-2m)_++\gamma}
      \nonumber\\
    &\le \frac{2m}{P(B)}+1+
      \frac{\gamma}{P(B)}\log\frac{((N-1)P(B)-2m)_++\gamma}{\gamma}
      \nonumber\\
    &\le \frac{1}{P(B)}\okra{2m+
      \frac{\gamma}{P(B)}\log\frac{N-1+\gamma}{\gamma}}+1
      .
  \end{align}
  On the other hand, almost surely, we have
  \begin{align}
    \min_{B\in\beta} C^B_m\le N:=(m-1)\card\beta+1
  \end{align}
  since in $(m-1)\card\beta+1$ consecutive draws from set $\beta$
  there must be an element $B\in\beta$ that is repeated $m$ times.
  Consequently, we may write
  \begin{align}
    \mean \min_{B\in\beta} C^B_m
    &\le \mean \min\klam{\min_{B\in\beta} C^B_m,N}
      \le \min_{B\in\beta} \mean \min\klam{C^B_m, N}
      .
  \end{align}
  Hence we obtain the claim.
\end{proof}

Notice that the repetition time can be expressed as
\begin{align}
  R^{(2)}_k=\min_{x_1^k} C^{(X_1^k=x_1^k)}_2.
\end{align}
Consequently, we can derive this upper bound for the repetition time.
\begin{proposition}
  \label{theoSandwichRepetitionOneOne}
  Consider a stationary process $(X_i)_{i\in\mathbb{Z}}$ over a
  $D$-ary alphabet.  Assume that $\sum_{k=1}^\infty \rho_k<\infty$ for
  some $\rho_k>0$.  Almost surely, for sufficiently large $k$, we have
  \begin{align}
    \label{SandwichRepetitionOneOne}
    \log R^{(2)}_k<H_\infty(X_1^k)+\log
    \okra{4+\gamma_k\log\frac{D^k+\gamma_k}{\gamma_k}}
    -\log\rho_k
    ,
  \end{align}
  where coefficients $\gamma_k$ are defined in (\ref{Gamma}).
\end{proposition}
\begin{proof}
  In view of the Markov inequality and Lemma \ref{theoShortMemory},
  we may bound
  \begin{align}
    P(R^{(2)}_k\ge n)
    &\le \frac{\mean R^{(2)}_k}{n}
      \le \frac{\mean \min_{x_1^k} C^{(X_1^k=x_1^k)}_2}{n}
      \nonumber\\
     &\le \frac{e^{H_\infty(X_1^k)}
      \okra{4+\gamma_k\log\frac{D^k+\gamma_k}{\gamma_k}}+1}{n}.
  \end{align}
  Define
  $n_k:=e^{H_\infty(X_1^k)}
  \okra{4+\gamma_k\log\frac{D^k+\gamma_k}{\gamma_k}}\rho_k^{-1}$. Then
  we have $\sum_{k=1}^\infty P(R^{(2)}_k\ge n_k)<\infty$. Hence
  inequality (\ref{SandwichRepetitionOneOne}) is obtained by the
  Borel-Cantelli lemma.
\end{proof}

Using similar techniques, we may demonstrate an upper bound for the
expectation of variable $\okra{\min_{B\in\beta} C^B_m}^{-1}$, where $\beta$
is a finite partition.
\begin{lemma}
  \label{theoMaxBounded}
  Assume a stationary automorphism $T$ and a finite partition $\beta$.
  Suppose that
  \begin{align}
    \delta:=
    \sup_{n\in\mathbb{N}}\frac{\mean\max_{B\in\beta}
    F^B_n}{\max_{B\in\beta}\mean F^B_n}<\infty.
  \end{align}
  Then for $m\in\mathbb{N}$, we have
  \begin{align}
    \mean \okra{\min_{B\in\beta} C^B_m}^{-1}\le
    \max_{B\in\beta} P(B)\, \frac{\delta}{m}\okra{
    \log\frac{m\card\beta}{\delta}+1
    }
    .
  \end{align}
\end{lemma}
\begin{proof}
  In view of Theorem \ref{theoExtreme}, the Markov inequality, and
  inequality $\max_{B\in\beta} P(B)\ge (\card\beta)^{-1}$, we obtain
  \begin{align}
    \mean \okra{\min_{B\in\beta} C^B_m}^{-1}
    &=
      \int_1^\infty P\okra{\min_{B\in\beta} C^B_m\le \floor{n}} \frac{dn}{n^2}
      \nonumber\\
    &=
      \int_1^\infty P\okra{\max_{B\in\beta} F^B_{\floor{n}}\ge m} \frac{dn}{n^2}
      \nonumber\\
    &\le \int_1^\infty \min\klam{1,\frac{\mean\max_{B\in\beta}
      F^B_{\floor{n}}}{m}} \frac{dn}{n^2}
      \nonumber\\
    &\le \int_1^\infty \min\klam{1,\frac{\delta n\max_{B\in\beta}
      P(B)}{m}} \frac{dn}{n^2}
      \nonumber\\
    &= \frac{\delta \max_{B\in\beta}
      P(B)}{m}\okra{\log\frac{m}{\delta \max_{B\in\beta}
      P(B)}+1}
      \nonumber\\
    &\le \max_{B\in\beta} P(B)\, \frac{\delta}{m}\okra{
      \log\frac{m\card\beta}{\delta}+1
      }
      .
  \end{align}
  Thus we obtain the claim.
\end{proof}

Consequently, let us derive a lower bound for the repetition time.
\begin{proposition}
  \label{theoSandwichRepetitionOneTwo}
  Consider a stationary process $(X_i)_{i\in\mathbb{Z}}$ over a
  $D$-ary alphabet.  Assume that $\sum_{k=1}^\infty \rho_k<\infty$ for
  some $\rho_k>0$.  Almost surely, for sufficiently large $k$, we have
  \begin{align}
    \label{SandwichRepetitionOneTwo}
    \log R^{(2)}_k>H_\infty(X_1^k)-\log \frac{\delta_k}{2}
    -\log\okra{\log\frac{2D^k}{\delta_k}+1}
    +\log\rho_k
    ,
  \end{align}
  where coefficients $\delta_k$ are defined in (\ref{Delta}).
\end{proposition}
\begin{proof}
  In view of the Markov inequality and Lemma \ref{theoMaxBounded},
  we may bound
  \begin{align}
    P(R^{(2)}_k\le n)
    &\le n\mean \okra{R^{(2)}_k}^{-1}
      \le n\mean \okra{\min_{x_1^k} C^{(X_1^k=x_1^k)}_2}^{-1}
      \nonumber\\
    &\le n e^{-H_\infty(X_1^k)}
      \frac{\delta_k}{2}\okra{\log\frac{2D^k}{\delta_k}+1}.
  \end{align}
  Define
  $n_k:=e^{H_\infty(X_1^k)}
  \frac{2}{\delta_k}\okra{\log\frac{2D^k}{\delta_k}+1}^{-1}\rho_k$. Then
  we have $\sum_{k=1}^\infty P(R^{(2)}_k\le n_k)<\infty$. Hence
  inequality (\ref{SandwichRepetitionOneTwo}) is obtained by the
  Borel-Cantelli lemma.
\end{proof}

\section{Yet another lower bound}
\label{secYetAnother}

We also have a more general lower bound for the repetition time.  The
following proposition refines \citep[Theorem 9.25]{Debowski21} by
using the weighted conditional entropy rather than the min-entropy
conditioned on the infinite future.
\begin{proposition}
  \label{theoSandwichRepetitionTwo}
  Consider a stationary process $(X_i)_{i\in\mathbb{Z}}$ over a
  countable alphabet. Assume that $\sum_{k=1}^\infty \rho_k<\infty$
  for some $\rho_k>0$. Define the weighted conditional entropy
  \begin{align}
    H_-(X_1^k|X_{k+1}^\infty):=-\log \okra{
    \sum_{i=1}^\infty\frac{e^{-H_\infty(X_1^k|X_{k+1}^{k+i})}}{i(i+1)}}.
  \end{align}
  Almost surely, for sufficiently large $k$, we have
  \begin{align}
    \label{SandwichRepetitionTwo}
    \log R^{(2)}_k> \frac{1}{3}H_-(X_1^k|X_{k+1}^\infty)
    +\frac{1}{3}\log\rho_k
    .
  \end{align}
\end{proposition}
\emph{Remark:} We note that
$H_-(X_1^k|X_{k+1}^\infty)\ge
H_\infty(X_1^k|X_{k+1}^\infty)$. Hence Proposition
\ref{theoSandwichRepetitionTwo} is stronger than \citep[Theorem
9.25]{Debowski21}.
\begin{proof}
  For $i<j$, we obtain 
\begin{align}
  P(X_{i+1}^{i+k}=X_{j+1}^{j+k})
  &=
    \mean
    P(X_{i+1}^{i+k}=X_{j+1}^{j+k}|X_{i+k+1}^{j+k})
    \nonumber\\
  &\le
    \mean \max_{x_1^k} P(X_{i+1}^{i+k}=x_1^k|X_{i+k+1}^{j+k})
    \nonumber\\
  &=
    e^{-H_\infty(X_1^k|X_{k+1}^{k+j-i})}
    .
\end{align}
Observe that $3-2x\le x^{-2}$, so $n-x\le 4n^3x^{-2}/27$.
%$n-x=n/3(3-2(3x/2n))\le n/3(3x/2n)^{-2}=4n^3/27x^2$.
Hence, by the union bound,
\begin{align}
  P(R^{(2)}_k\le n)
  &= P\okra{X_{i+1}^{i+k}=X_{j+1}^{j+k} \text{ for some }
    0\le i<j\le n}
    \nonumber\\
  &\le \sum_{0\le i<j\le n} P(X_{i+1}^{i+k}=X_{j+1}^{j+k})
    \le \sum_{0\le i<j\le n}  e^{-H_\infty(X_1^k|X_{k+1}^{k+j-i})}
    \nonumber\\
  &= \sum_{i=1}^{n-1}
    \frac{n-i}{e^{H_\infty(X_1^k|X_{k+1}^{k+i})}}
    \le \frac{8n^3}{27}
    \sum_{i=1}^\infty\frac{e^{-H_\infty(X_1^k|X_{k+1}^{k+i})}}{i(i+1)}
    \nonumber\\
  &\le \frac{8n^3}{27} e^{-H_-(X_1^k|X_{k+1}^\infty)}.
\end{align}
Put $n_k:=\sqrt[3]{\rho_k e^{H_-(X_1^k|X_{k+1}^\infty)}}$. Then
$\sum_{k=1}^\infty P(R^{(2)}_k\le n_k)<\infty$.  Consequently,
inequality (\ref{SandwichRepetitionTwo}) follows by the
Borel-Cantelli lemma.
\end{proof}

The last proposition in this section recycles some techniques used in
universal coding such as \citep[Theorem 6.11]{Debowski21}. The goal is
to show that the bound in terms of the weighted conditional entropy is
more precise than the bound proved in \citep[Theorem 8]{Debowski18b}
since the length of conditioning can be shortened to $I_k$ rather than
$D^k$.
\begin{proposition}
  \label{theoSandwichContext}
  Consider a stationary process $(X_i)_{i\in\mathbb{Z}}$ over a
  countable alphabet.  For $I_k$ defined in (\ref{ContextLength}), we
  have
\begin{align}
  \label{SandwichContextOne}
  \log I_k-\log 2&\le H_-(X_1^k|X_{k+1}^\infty)\le
  3\log I_k+\frac{1}{I_k}
  ,
  \\
  \label{SandwichContextThree}
  \log (I_k-1)-H_0(X_i)&\le H_\infty(X_1^k|X_{k+1}^{k+I_k})\le \log I_k
                              .
\end{align}
\end{proposition}
\emph{Remark:} We note that $I_k\le D^k$ and
$H_\infty(X_1^k|X_{k+1}^{k+I_k})\ge H_\infty(X_1^k|X_{k+1}^{k+D^k})$
for a $D$-ary alphabet. Hence Propositions
\ref{theoSandwichRepetitionTwo} and \ref{theoSandwichContext} combined
are stronger than \citep[Theorem 8]{Debowski18b}.
\begin{proof}
First, we derive (\ref{SandwichContextOne}) from
\begin{align}
  H_-(X_1^k|X_{k+1}^\infty)
  &\le
    H_\infty(X_1^k|X_{k+1}^{k+I_k})+\log I_k+\log (I_k+1)\le
    3\log I_k+\frac{1}{I_k},
  \\
  H_-(X_1^k|X_{k+1}^\infty)
  &\ge
    -\log\okra{\okra{1-\frac{1}{I_k}}e^{-H_\infty(X_1^k|X_{k+1}^{k+I_k-1})}
    +\frac{1}{I_k}}
    \nonumber\\
  &\ge -\log\okra{\okra{1-\frac{1}{I_k}}\frac{1}{I_k-1}+\frac{1}{I_k}}
    =\log I_k-\log 2.
\end{align}
Moreover, for a stationary process,
\begin{align}
  0
  &\le
    H_\infty(X_1^k|X_{k+1}^{k+I_k-1})-H_\infty(X_1^k|X_{k+1}^{k+I_k})\le
    H_0(X_i)
\end{align}
since by the chain rule (\ref{Chain}), we have
\begin{align}
  H_\infty(Y,X|Z,V)\le H_\infty(Y,X|Z)
  &\le H_0(X|Z)+H_\infty(Y|X,Z)
    \nonumber\\
  &\le H_0(X)+H_\infty(U,Y|X,Z). 
\end{align}
Hence we obtain (\ref{SandwichContextThree}).
\end{proof}

\section{Final touch}
\label{secFinal}

Let us recall some properties of Hilberg exponents.  The probabilistic
results of \citep[Section 8.1]{Debowski21} can be summarized as
follows.
\begin{definition}
  Let $T$ be an automorphism of the probability space.  A sequence of
  random variables $(J_k)_{k\in\mathbb{N}}$ is called $T$-increasing
  if
  \begin{align}
    J_{k+1}\ge J_k\circ T,J_k\ge 0.
  \end{align}
\end{definition}
Let $\median X:=\sup \klam{r: P(X<r)\le \frac{1}{2}}$ denote the
median of a variable $X$.
\begin{lemma}
  \label{theoHilberg}
  For a stationary ergodic automorphism $T$ and a $T$-increasing
  sequence $(J_k)_{k\in\mathbb{N}}$, we have
  \begin{align}
    \hilberg_{k\rightarrow\infty} \median J_k
    \le \hilberg_{k\rightarrow\infty} J_k 
    \le \hilberg_{k\rightarrow\infty} \mean J_k
    \text{ a.s.}, 
  \end{align}
  where random variable $\hilberg_{k\rightarrow\infty} J_k$ is almost
  surely constant. Moreover, the inequalities become the equalities if
  \begin{align}
    \limsup_{k\rightarrow\infty} \frac{\sqrt{\var J_k}}{\mean J_k}<1.
  \end{align}
\end{lemma}
\begin{proof}
  All claims but the last one follow by \citep[Theorems 8.3--8.5]{Debowski21}.
  The last claim is a slight strengthening of \citep[Theorem
  8.6]{Debowski21} and in a similar fashion it follows by the
  median-mean inequality $\abs{\median X-\mean X}\le \sqrt{\var X}$,
  see e.g.\ \citep[Theorem 3.17]{Debowski21}.
\end{proof}

We also make this simple observation.
\begin{proposition}
  \label{theoIncreasing}
  Consider the shift automorphism $X_i\circ T=X_{i+1}$. Then sequences
  $(L^{(\gamma)}_n)_{n\in\mathbb{N}}$,
  $(R^{(\gamma)}_k)_{k\in\mathbb{N}}$,
  $(-\log P(X_1^k))_{k\in\mathbb{N}}$, and
  $(-\log P(X_1^k|X_{k+1}^\infty))_{k\in\mathbb{N}}$ are
  $T$-increasing.
\end{proposition}
Thus Theorem \ref{theoMainOne} follows by Propositions
\ref{theoSandwichRecurrence} and \ref{theoIncreasing} and Lemma
\ref{theoHilberg}, whereas Theorem \ref{theoMainTwo} follows by
Propositions
\ref{theoSandwichRepetitionOneOne}--\ref{theoSandwichRepetitionOneTwo}
and Theorem \ref{theoMainThree} follows by Propositions
\ref{theoSandwichRepetitionTwo}--\ref{theoSandwichContext}. Quod erat
demonstrandum.

% \bibliographystyle{IEEEtran}
\bibliographystyle{abbrvnat}
%\bibliographystyle{elsart-harv}

\bibliography{0-journals-abbrv,0-publishers-abbrv,ai,mine,tcs,ql,nlp,books}

\end{document}

%%% Local Variables: 
%%% mode: latex
%%% mode: flyspell
%%% TeX-master: t
%%% TeX-PDF-mode: t
%%% ispell-local-dictionary: "american"
%%% End: 

%  LocalWords:  differentiable
